\subsection{Global sections realizing each finite sharp theta jet}

\begin{theorem}[DP54: finite-jet globalization at square tensor levels]\label{DP54}
Use the actual polarized Jacobian of DP49 over
$R=W(\mathbf F_{11^3})$, with fraction field $K$, and put
\[
 \mathcal L=(1,w_0)^*\mathcal P_\lambda,
 \qquad w_0=-(a+a^2+a^4).
\]
This trace-one theta line is relatively ample.
In the invariant product canonical gauge of DP47, fix an integral formal frame and write the normalized germ of a section $s$ of $\mathcal L^{\otimes\kappa}$ as
\[
 \Theta_s(z)=f_s(z)U(z)^\kappa B(z)^{-\kappa},
 \qquad U\in1+(z)R[[z]],\qquad B=\exp b.
\]
Here $B$ is the actual primitive arbitrary-unit Poincar\'e splitting of DP46.
Let $E$ be the inverse normalized group logarithm and put $\beta=b\circ E$.

For every $n\ge0$ there is $\kappa_0(n)$ such that for every integer
$\kappa\ge\kappa_0(n)$ some global integral section
$s_\kappa\in H^0(J,\mathcal L^{\otimes\kappa})$ satisfies
\begin{equation}\label{eq:globalfinitejet}
 \Theta_{s_\kappa}(E(Z))
 \equiv\exp(-\kappa\beta(Z))\pmod{(Z)^{n+1}}.
\end{equation}
Its algebraic coefficient is a unit at zero. The same section used in any fixed finite number of copies has the corresponding product pure-kernel jet through total degree $n$.

Fix finitely many linear functionals on finite jets of those copies whose pure-kernel values at exponent $-1$ are nonzero. For some integer $B_0$, every sufficiently large positive integer
\[
 M\equiv1\pmod{11^{B_0}},\qquad \kappa=M^2,
\]
admits such a section for which every selected value has exactly the valuation of its pure-kernel value at exponent $-1$.
For a homogeneous degree-$n$ differential functional, multiplication by the source factor $M^{-n}$ preserves that valuation.

Consequently the sharp ordinary primitive-line sequences of DP46, and the first-negative-row two-column harmonic sequence of DP51, are realized in valuation by global sections of varying square tensor powers of this actual theta line. Their rates are respectively $1/1220$ or $1/120$ for the two primitive-line classes, and $1/1220$ for the first-row harmonic family.
The sections may also be chosen over $\Q(\zeta_7)$, integral at the selected prime above $11$, preserving these finite valuations.
The section and tensor power may vary with the derivative order; no tensor-size, height, cost, or prescribed translated-family bound is asserted.
\end{theorem}

\begin{proof}
\emph{Relative ampleness with the correct sign.}
On the geometric closed fibre let $\mathcal T$ be an ample theta line inducing the canonical principal polarization. DP49 proves
\[
 \phi_{\mathcal L}=\phi_{\mathcal T}=\lambda.
\]
Thus $\mathcal N=\mathcal L\otimes\mathcal T^{-1}$ has
$\phi_{\mathcal N}=0$ and belongs to $\operatorname{Pic}^0$.
The known Picard facts needed here are Milne, \emph{Abelian Varieties} (1986), Proposition~9.2, pp.~117--118, the following definition of $\operatorname{Pic}^0$ on p.~118, and Proposition~10.1, p.~119: an ample line over an algebraically closed field has a surjective map $\phi_{\mathcal T}$ onto $\operatorname{Pic}^0$.
Choose $c$ such that
\[
 \mathcal N\cong t_c^*\mathcal T\otimes\mathcal T^{-1}.
\]
Then $\mathcal L\cong t_c^*\mathcal T$ is ample on the geometric closed fibre.
Ampleness of the fixed line descends along the ground-field extension, by Stacks Project Lemma~35.25.6, and closed-fibre ampleness spreads to an open neighbourhood by Lemma~30.21.4.
Every open neighbourhood of the closed point of $\operatorname{Spec}R$ is the whole spectrum. Therefore $\mathcal L$ is relatively ample, and is ample on $J$ because the base is affine. This uses equality with the positive principal polarization, not merely that $\phi_{\mathcal L}$ is an isogeny.\footnote{The exact primary locators are J.~S.~Milne, ``Abelian Varieties,'' in \emph{Arithmetic Geometry} (1986), pp.~117--119, official author copy \url{https://www.jmilne.org/math/articles/1986b.pdf}; Stacks Project Lemma~35.25.6 in \url{https://stacks.math.columbia.edu/tag/02LB}, and Lemma~30.21.4, \url{https://stacks.math.columbia.edu/tag/0D2N}.}

\emph{Integral finite-jet extension.}
Let $I$ be the ideal of the zero section and $Z_n=V(I^{n+1})$.
Smoothness and the formal parameters give
\[
 \mathcal O(Z_n)
 =R[[z_1,z_2,z_3]]/(z_1,z_2,z_3)^{n+1}.
\]
This is finite free over $R$. The formal frame restricts to a genuine frame on this finite thickening; no algebraization of the entire infinite frame is required.
The ideal sheaf $I^{n+1}$ is coherent, so Serre vanishing gives
\[
 H^1(J,I^{n+1}\otimes\mathcal L^{\otimes\kappa})=0
 \qquad(\kappa\ge\kappa_0(n)).
\]
The exact known form is Stacks Project Lemma~30.16.1(4), or its relative version Lemma~30.16.2.\footnote{\url{https://stacks.math.columbia.edu/tag/0B5T}, part~(4), and \url{https://stacks.math.columbia.edu/tag/02O1}.}
The restriction sequence therefore yields the integral surjection
\[
 H^0(J,\mathcal L^{\otimes\kappa})
 \twoheadrightarrow
 H^0(Z_n,\mathcal L^{\otimes\kappa}|_{Z_n}).
\]
Prescribe $f_s\equiv U^{-\kappa}\pmod{I^{n+1}}$ in the restricted frame and lift it. This target is integral and has value one at zero. Thus $s$ is nonzero and
$f_sU^\kappa\equiv1\pmod{I^{n+1}}$.
Substitution by $E$ preserves total order because it has zero constant term and identity tangent map. Multiplication by $\exp(-\kappa\beta)$ also preserves the ideal of terms of degree at least $n+1$. This proves \eqref{eq:globalfinitejet}.
In a product of finitely many copies, every nonconstant prefactor error already has total degree at least $n+1$, so the same section realizes the product jet.

\emph{Polynomial dependence on the tensor exponent.}
For a fixed finite jet functional $\ell$ of order at most $n$ on $r$ copies, put
\[
 F_\ell(T)=\ell\left(
 \exp\left[-T\sum_{b=1}^r\beta(Z^{(b)})\right]\right).
\]
This lies in $K[T]$: the logarithmic sum has zero constant term, so only its powers of order at most $n$ can contribute to the chosen jet.
Let $m_\ell$ be the minimum valuation of the nonzero coefficients of this polynomial. For an integer $T\equiv1\pmod{11^B}$,
\[
 v_{11}\bigl(F_\ell(T)-F_\ell(1)\bigr)\ge B+m_\ell,
\]
since $T^a-1$ is divisible by $11^B$ for every positive integer $a$.
Choose $B$ with
\[
 B+m_\ell>v_{11}(F_\ell(1)).
\]
Then $v_{11}(F_\ell(T))=v_{11}(F_\ell(1))$.
For finitely many selected nonzero values take the maximum of the finitely many required bounds.
There are arbitrarily large positive integers $M\equiv1\pmod{11^B}$; choose any with $M^2\ge\kappa_0(n)$ and set $T=M^2$.
The exact finite-jet lift then proves all the selected valuation equalities.
The source's definition \texttt{thExpr}, \texttt{IsThRule.lean}:66--70, uses tensor exponent $M^2$ and factor $M^{-\deg P}$. Here the latter is a unit, so it does not change the valuation.

\emph{Sharp sequences.}
For a primitive logarithmic line of the actual pair in DP49, DP46 gives
\[
 \liminf_{n\longrightarrow\infty}
 \frac1n v_{11}\bigl(n![t^n]\exp[-\beta(ut,vt)]\bigr)
 =c_k,
\]
where $c_2=1/1220$ if the first dual component is a unit and $c_1=1/120$ if the first is divisible by $11$ and the second is a unit.
Choose a sequence of nonzero values attaining this liminf and apply the finite-jet argument separately at each selected degree.
It produces global sections $s_{M_\nu^2}$ with exactly the same selected valuations, including the unit factor $M_\nu^{-n_\nu}$.

For the first dual row the two-column integral polynomials
\[
 P_{a,b}=\sum_{i=0}^{\min(a,b)}(-1)^i
 \binom ai\binom bi
 X_1^{a-i}Y_{1,1}^{b-i}X_2^iY_{1,2}^i
\]
are the source-shaped harmonic family of DP51.
Its proof applies to the actual arbitrary-unit display: DP46 retains the same unique initial monomial with unit coefficient, the gap to the other terms, and the primitive-path threshold used in that proof.
It supplies nonzero values of rate $1/1220$ as $a+b$ tends to infinity. Use two copies of the same finite-jet lift with $n=a+b$, and apply the same polynomial congruence argument in the tensor exponent. The resulting product of global normalized sections realizes each selected valuation.
This gives a sequence of sections and tensor powers, not the asymptotic sequence of a fixed section or a fixed power.

\emph{Number-field sections.}
Let $E_0=\Q(\zeta_7)$, choose the prime $\mathfrak p\mid11$ underlying the given embedding, and let $R_0=\mathcal O_{E_0,\mathfrak p}$. Its completion is $R$.
The curve, its Jacobian, $a$, $w_0$, the polarization, and $\mathcal L$ are defined over $R_0$ by the good model and formulas of DP49.
For a fixed tensor power already chosen, the module
\[
 V_0=H^0(J_{R_0},\mathcal L^{\otimes\kappa})
\]
is finite and torsion-free over the DVR $R_0$, hence free. Proper coherent cohomology gives finiteness; flatness of $J_{R_0}$ and invertibility of the sheaf give injectivity of multiplication by $11$, hence torsion-freeness.
Flat cohomology base change gives
\[
 V_0\otimes_{R_0}R
 =H^0(J_R,\mathcal L^{\otimes\kappa}).
\]
The exact locator is Stacks Project Lemma~30.5.2(2), Tag~02KH.\footnote{\url{https://stacks.math.columbia.edu/tag/02KH}.}
Density of $R_0$ in its completion makes $V_0$ dense in the right-hand side in a finite free basis. Each selected finite jet evaluation is a finite linear expression in the section coordinates; for a product of copies it is a finite polynomial expression. It is therefore continuous. A sufficiently close section in $V_0$ preserves every prescribed nonzero valuation and its unit value at zero.
This supplies a number-field section integral at $\mathfrak p$, with the selected valuations preserved. It need not preserve the exact formal coefficients, and it gives no height bound.
\end{proof}

The theorem removes the objection that the sharp prefactor-$1$ examples exist only as formal germs: each relevant finite jet is realized by global algebraic theta sections at sufficiently large square tensor levels. What remains is quantitative and specific to the prescribed family: bounds for the tensor level, the section heights and norms, and the realization inside the required translated family with its endpoint and nonvanishing inputs. No such bound, family membership, or shorter chain follows from this existence argument.
